% !TEX root = ../main.tex
\section{Results and Discussion (Preliminary)}
\label{sec:results}

% Revised 2026-10-11 from the owner's hand review of the printed report
% (handwritten/2026-10-11_main.md, notes 17.x to 21.x). Renamed from
% 04_results_discussion.tex. Removed as struck: both scope notes and the
% Pareto-front call-out. Moved out: the options table, the envelope figure and
% the "Limits" call-out (to Sections 3 and 2). Moved in: the emissions
% subsection from the Impact section. Passages between \aidraft{...} and
% "% ---- end AI draft" were drafted by Claude and are the owner's to rewrite.

\irtwo{Results and Discussion (preliminary). For Interim Report \#2 this section needs: (1) the architecture as it stands (context diagram, aircraft IBD, the three requirements tagged \texttt{improved} as the measured gap) --- the figures are now in Sections~\ref{sec:scope} and~\ref{sec:architecture-results}; (2) the two hand checks from \texttt{experiment-definition.md} \S9 written up as preliminary results; (3) the first OpenAP numbers --- Figure~\ref{fig:fuelflow-altitude} and Table~\ref{tab:outcomes} are in place, the discussion is not; (4) one paragraph on what is not yet done.}

\subsection{Operational Concept}
\label{sec:conops-results}

% Owner's notes, 2026-10-11 print, page 17: "DO"; "Show full detail 1st, then w/
% abstractions & provide rationale."; "The conops doesn't present my specific
% scenario. Need to update w/ the example conflict."
\aidraft{The operational concept: full detail, then the abstraction and its reason, then the example conflict (your notes: ``DO''; ``Show full detail 1st, then w/ abstractions and provide rationale''; ``Need to update w/ the example conflict''). Sources: \texttt{knowledge/models/conops-nominal-domestic\_flight.md}, \texttt{knowledge/models/conops-hub-to-hub-trajectory-cost.md} (both marked as working documents, not reconciled against the literature), D-016 item 5, and the activity diagram. \texttt{FlyEnroute} is not yet decomposed in the model, so the last paragraph describes where the conflict will sit, not what the diagram shows today.}
The operational concept was first written in full detail as a gate-to-gate thread for one nominal commercial flight, from the pilot's point of view. It runs through eleven phases: the airline's schedule, fleet and crew assignment; the pilot's preparation before duty; reporting and the dispatch release; aircraft preflight and initialization of the flight management system; boarding, loading and fueling; clearance, pushback and taxi; takeoff and initial climb; climb; cruise; arrival planning; and descent through to the gate. Each phase records its initial state, the interfaces in use, and the questions it raised. The full thread is in Appendix~\ref{app:conops}.

The model carries an abstraction of that thread. Figure~\ref{fig:sysml-nominalOpsActivity} shows six actions in sequence: plan the flight, prepare for departure, depart, fly en route, arrive, and close out the flight. The eleven phases were reduced to six because the study concerns one decision in cruise. The phases before departure and after arrival fix the conditions that decision inherits (the aircraft's weight, its cost index, its schedule position) and are kept at the level needed to show where those conditions come from. Detail is added only where the studied decision is made.

The example conflict sits inside the fly en route action. Two aircraft in cruise are predicted to cross at the same level at the same time. The sector controller must move one of them, and each choice costs fuel, time and instructions that fall on different parties. The nominal thread treats a trajectory change in cruise as one step; the experiment of Section~\ref{sec:experiment} opens that step and asks what the decider knows when making it.
% ---- end AI draft

\plantodo{S5}{Decompose \texttt{FlyEnroute} one level in \texttt{nas\_activity\_diagram-nominal\_ops.sysml} so the conflict-resolution step appears on Figure~\ref{fig:sysml-nominalOpsActivity} (D-016 item 5), then check the three paragraphs above against the diagram. Put the full activity breakdown in Appendix~\ref{app:conops}.}

% GENERATED by tools/render_sysml_diagrams.py. Do not edit: the next run overwrites it.
% To change a caption, edit the `// caption:` line above the view in the .sysml file.
% To add or remove a figure, add or remove the view in the model.
% To move a figure between the body and the appendix, add or remove a
% `// placement: body` line above the caption line.

\begin{landscape}
\begin{figure}[!htbp]
\centering
\pdfimageresolution=300
\includegraphics[width=\linewidth,height=0.8\textheight,keepaspectratio]{sysml/diagram-nominalOpsActivity}
\caption{Nominal flight operations activity diagram, phase level (each phase is decomposed in its own figure). View \texttt{nominalOpsActivity}; source \texttt{nas\_activity\_diagram-nominal\_ops.sysml}; rendered 2026-10-18.}
\label{fig:sysml-nominalOpsActivity}
\end{figure}
\end{landscape}


\subsection{Stakeholder Objectives}
\label{sec:objectives-results}

% Owner's note, 2026-10-11 print, page 17: "DO".
\aidraft{What ``optimal'' means to each stakeholder, and where the objectives conflict (your note: ``DO''). Source: the synthesis section of \texttt{knowledge/models/stakeholder-objective-ontology.md}, whose own status note says the classifications are judgment calls not yet sourced from annotated literature. Limited here to the stakeholders and objectives kept in Section~\ref{sec:scope}.}
Each stakeholder kept in scope pursues a different objective when an aircraft's trajectory is changed. The airline wants the lowest cost for its flights, which combines fuel, time, and the consequences of a late arrival for connecting passengers. Air traffic control wants separation assured with a workload the sector team can sustain, and a traffic picture that stays predictable. The flight crew wants instructions the aircraft can fly, given with enough time to set them up.

Some of these objectives align. Burning less fuel lowers the airline's cost and lowers carbon dioxide emissions together. Air traffic control's wish for predictability and the airline's wish to keep its schedule both favor aircraft that fly what was planned.

Others conflict, and the conflict this project studies is the first of them: the trajectory that is cheapest for one aircraft is not always the one that costs the sector least to grant. A request or a resolution that saves fuel for one flight can add instructions, monitoring, and follow-up for the controller. That cost is externalized. It does not appear in the airline's own cost function, and the party that bears it has no means to price it back.

The objectives also operate on different timescales. An airline sets its schedule policy over months, realizes fuel cost per flight, and realizes delay cost in minutes. Air traffic control manages capacity over years through sector design and over seconds through tactical control. A decision made in minutes by one party is judged against objectives the others hold on other clocks.
% ---- end AI draft

\plantodo{S9 (owner's note, 2026-10-11 print: ``Not really here yet. Can show prelim results from my table'')}{The myopic-optimization analysis (definition of a local optimum in the system-of-systems setting, candidate system-level measures of value, the documented conflicts) is not written. For this report, show it through the one conflict the experiment measures: read across the rows of Table~\ref{tab:outcomes} for an option that is better for one aircraft or airline and worse for the sector or the system. The other candidate conflicts are future work (D-016) and go to Appendix~\ref{app:myopic}.}

\subsection{Architecture: Detailed and Abstracted}
\label{sec:architecture-results}

% Owner's note, 2026-10-11 print, page 17: "Place detailed & abstracted. Make sub-section."
\plantodo{S10}{Present the SysML architecture: package structure, the context diagram (Figure~\ref{fig:sysml-systemContext}), the internal block diagrams below, the interface and item-flow model, and the requirements model with traceability. For the aircraft, the detailed diagram comes first and the abstracted one second (your note: ``Place detailed and abstracted''); say what the abstraction drops and why. The airspace management diagram is still a placeholder.}

\aidraft{Two connecting sentences so the figures are referred to from the text. Replace with your own description of what each level of detail shows.}
The aircraft is shown at two levels of detail. Figure~\ref{fig:sysml-aircraftIbdDetail} is the detailed internal block diagram, including propulsion, airframe, landing gear, utilities, and the fuel and power links. Figure~\ref{fig:sysml-aircraftIbd} is the abstraction used in the rest of the report: the crew, the flight deck controls, the avionics, the flight controls, the fuel system, and the information and command links among them, which are the parts the studied decision passes through. Figure~\ref{fig:sysml-airspaceManagementIbd} shows the internals of airspace management as modeled so far.
% ---- end AI draft

% GENERATED by tools/render_sysml_diagrams.py. Do not edit: the next run overwrites it.
% To change a caption, edit the `// caption:` line above the view in the .sysml file.
% To add or remove a figure, add or remove the view in the model.
% To move a figure between the body and the appendix, add or remove a
% `// placement: body` line above the caption line.

\begin{landscape}
\begin{figure}[!htbp]
\centering
\pdfimageresolution=300
\includegraphics[width=\linewidth,height=0.8\textheight,keepaspectratio]{sysml/diagram-aircraftIbdDetail}
\caption{Aircraft internal block diagram, detailed level: adds propulsion, airframe, landing gear, utilities, and the fuel and power links. View \texttt{aircraftIbdDetail}; source \texttt{nas\_ibd\_aircraft.sysml}; rendered 2026-10-19.}
\label{fig:sysml-aircraftIbdDetail}
\end{figure}
\end{landscape}

% GENERATED by tools/render_sysml_diagrams.py. Do not edit: the next run overwrites it.
% To change a caption, edit the `// caption:` line above the view in the .sysml file.
% To add or remove a figure, add or remove the view in the model.
% To move a figure between the body and the appendix, add or remove a
% `// placement: body` line above the caption line.

\begin{landscape}
\begin{figure}[!htbp]
\centering
\pdfimageresolution=300
\includegraphics[width=\linewidth,height=0.8\textheight,keepaspectratio]{sysml/diagram-aircraftIbd}
\caption{Aircraft internal block diagram, default level of detail. View \texttt{aircraftIbd}; source \texttt{nas\_ibd\_aircraft.sysml}; rendered 2026-10-19.}
\label{fig:sysml-aircraftIbd}
\end{figure}
\end{landscape}

% GENERATED by tools/render_sysml_diagrams.py. Do not edit: the next run overwrites it.
% To change a caption, edit the `// caption:` line above the view in the .sysml file.
% To add or remove a figure, add or remove the view in the model.
% To move a figure between the body and the appendix, add or remove a
% `// placement: body` line above the caption line.

\begin{landscape}
\begin{figure}[!htbp]
\centering
\pdfimageresolution=300
\includegraphics[width=\linewidth,height=0.8\textheight,keepaspectratio]{sysml/diagram-airspaceManagementIbd}
\caption{Airspace management internal block diagram. View \texttt{airspaceManagementIbd}; source \texttt{nas\_ibd\_airspace\_management.sysml}; rendered 2026-10-10.}
\label{fig:sysml-airspaceManagementIbd}
\end{figure}
\end{landscape}


\subsection{Experiment Results}
\label{sec:experiment-results}

% The table body is generated by simulation/experiment_scoring.py (do not edit the
% file it inputs). It replaces the hand-typed table whose cells were empty.
\begin{table}[h]
\centering
\small
\caption{Outcome of each case at four levels, for the base scenario scored at the actual masses. Each cost is a difference from flying on undisturbed. Preliminary: the scenario's numbers are not all confirmed.}
\label{tab:outcomes}
% Generated by simulation/experiment_scoring.py; do not edit. Fidelity: actualMass
% Scenario: base; run 2026-10-29
\begin{tabular}{@{}llp{3.0cm}p{3.2cm}p{2.6cm}p{2.4cm}@{}}
\toprule
Case & Option & Largest cost to any one aircraft (kg fuel-equiv.) & Airline totals (USD) & Sector (calls / effort / risk) & System fuel, CO2 (kg) \\
\midrule
A & O10 & 40 (AC2) & X 0, Y 102, Z 0 & 1+1 / 1 / low & 41, 130 \\
B & O8 & 26 (AC1) & X 66, Y 0, Z 0 & 1+1 / 3 / low & 32, 102 \\
C & O4 & 12 (AC2) & X 0, Y 32, Z 0 & 1+1 / 3 / low & 12, 39 \\
D & O4 & 12 (AC2) & X 0, Y 32, Z 0 & 1+1 / 3 / low & 12, 39 \\
\bottomrule
\end{tabular}

\end{table}

\plantodo{S12}{\textbf{Which option each case chose, and why.} Walk Table~\ref{tab:outcomes} row by row. For case A, name the rule that produced the choice. For B, C and D, name the cost that decided it.}

\plantodo{S12}{\textbf{Did the climb rule matter?} With the crossing pair above FL350, cases A and B could not offer a climb and cases C and D could. Say whether case C used one, for which aircraft, and what it saved. If C did not climb either, say so: that is a result too.}

\plantodo{S12}{\textbf{B minus A, and C minus B.} The first is what automation is worth; the second is what the information exchange is worth. Compare the size of each with the prior figure of tens of gallons per flight (\cite{engility2014tasarAlaska}, p.~12) and say whether the difference is larger than the truth model's uncertainty.}

\plantodo{S9, S12}{\textbf{Does the best option change with the level?} Read across each row of Table~\ref{tab:outcomes}. Name any option that is better for one aircraft or one airline and worse for the sector or the system, and who pays (the largest cost to any one aircraft). This is the sub-system versus system discussion.}

\plantodo{S12}{\textbf{Airline X's two flights.} Discuss the option that touches both of one airline's aircraft: what the airline would choose between its late flight and its early one, and that the controller cannot see that preference. Connect to the dispatcher priorities (schedule integrity and connections rated at or above fuel; source of that table still to confirm).}

\plantodo{S12}{\textbf{Workload is three things.} Discuss clearances issued, effort to devise, and residual risk separately. Say where a ``more workload'' option on one of them is ``less'' on another.}

\plantodo{S12}{\textbf{When would the controller say no?} (owner's note, second hand review, 5.2). Discuss whether refusals are more likely in morning rushes at hub airports or near holidays, using the sector-loading rule in \cite{engility2014tasarAlaska} (p.~10). A hypothesis; needs a source.}

\plantodo{S12}{\textbf{Is schedule the right priority?} (owner's note, second hand review, 4.11). Contrast a carbon price, which makes fuel weigh more against time, with controllers' usual practice of giving schedule the priority.}

% ---- Figure slot added 2026-10-10 (AI-generated analysis figure). The file is
% written by simulation/b739_envelope.py (D-017).
\begin{figure}[h]
\centering
\figorbox{b739_fuelflow_altitude.png}{PLACEHOLDER: run \texttt{simulation/b739\_envelope.py} to produce \texttt{report/figures/b739\_fuelflow\_altitude.png}}
\caption{Fuel flow and specific range against altitude at Mach 0.80 for the heavy and light B737-9 (OpenAP, D-017).}
\label{fig:fuelflow-altitude}
\end{figure}

\plantodo{Track A}{\textbf{Figure slot: sequence diagram} of case C's conflict resolution (controller, ATC decision support, airline operations center, FMS, crew), rendered from the model once drafted (turn 3).}

\plantodo{S12}{\textbf{The case where nothing changes.} Report the variation with similar weights and the climb closed, where case C picks the same option as case B, to show the comparison is not rigged.}

\subsection{Trace Back to the Model}
\label{sec:traceback}

% Owner's note, 2026-10-11 print, page 19: "DO".
\aidraft{Where each column of the outcome table lives in the model (your note: ``DO''). Part and port names are from \texttt{cameo\_models/nas\_context\_diagram.sysml}; the requirement and its tag are from \texttt{knowledge/models/requirements-and-traceability.md}. Check the part chosen for each column.}
Each column of Table~\ref{tab:outcomes} belongs to a part of the context diagram (Figure~\ref{fig:sysml-systemContext}). The aircraft column belongs to the aircraft and its flight crew. The airline column belongs to the airline operations center within flight operations. The sector column belongs to air traffic control within airspace management. The system column belongs to no single part. It is a property of the NAS as a whole, which is the reason no constituent system computes it today and the reason the model is needed to state it.

The exchange that cases C and D add is one interface in the model: the airline operations center's performance share to air traffic control's decision support, carrying gross weight, cost index and the flight's delay margin. It satisfies system requirement SLR-INF-02, which is tagged as part of the improved architecture because today's NAS does not meet it. The interface is not drawn on Figure~\ref{fig:sysml-systemContext}, which shows the baseline.
% ---- end AI draft

\subsection{Emissions and Who Pays}
\label{sec:emissions-impact}

% Moved here from the Impact section per the owner's note on the 2026-10-11
% print, page 21: "This doesn't warrant its own section. Embed w/in Results
% discussion." The system-of-systems call-out the owner struck is removed.
\plantodo{S9, S12}{\textbf{Airline level: put CO2 into dollars.} Take the fuel difference between cases from Table~\ref{tab:outcomes}, convert to kg of CO2, then to dollars. State the route used and its source. Candidates: (1) the price of a carbon credit or allowance, if a scheme applies to the flight; (2) passenger value, meaning revenue gained or lost because some passengers choose the flight listed with lower emissions. Make the point that fuel is already counted once as fuel cost and that pricing its CO2 counts the same kilograms again, on purpose. SOURCES NOT YET HELD.}

\plantodo{S9, S12}{\textbf{Passenger value.} Discuss how passengers see a flight's emissions when booking, and whether that reaches the airline as demand. The author's own practice of choosing lower-emission listings is the motivation, not the evidence: cite a study. Note the gap between one decision in one sector and the label a passenger saw at booking; the link is the airline's average over many flights.}

\plantodo{S12}{\textbf{Scale it honestly.} One conflict in one sector is a small number. If a per-year or fleet-wide figure is offered, state the multiplier and its source (for example, how often level crossing conflicts occur), and label it illustrative.}

\plantodo{S9}{\textbf{Who pays.} For the option the system prefers, name the aircraft or airline that is worse off and by how much. This is the equity discussion.}
